\section{Satellite Capacity Estimation: Protocol and Data}
\label{app:volume}

This appendix records the operational detail behind the opening of
Section~\ref{sec:res_volume}, so that its measurements can be audited or repeated.

\subsection*{Data Sources}

All sources are public French open-data services, accessed over standard web map and feature
services and cached on disk so that a re-run does not depend on service availability.
Table~\ref{tab:app_volume_sources} lists each one with the role it plays in the pipeline.

\begin{table}[htbp]
\centering
\caption{Data sources of the capacity pipeline}
\label{tab:app_volume_sources}
\small
\begin{tabularx}{\linewidth}{>{\raggedright\arraybackslash}p{4.0cm} X}
\toprule
\textbf{Source} & \textbf{Use} \\
\midrule
Orthophotography, 20~cm & Primary imagery for detection and segmentation \\
Orthophotography, 5 to 10~cm & Partial coverage, used only for the resolution control of
Figure~\ref{fig:volume_resolution} \\
LiDAR surface model & Building height where available, delivered as 32-bit float raster \\
National topographic database & Building footprints and, where populated, height \\
Classified installations register & Declared storage volume, used only to calibrate the
fill factor \\
\bottomrule
\end{tabularx}
\end{table}

All geometry is handled in the national metric projection, in which the ground sample
distance is exactly the bounding-box width in metres divided by the image width in pixels
rather than an approximation.

\subsection*{Calibration of the Fill Factor}

The fill factor relates declared storage volume to the geometric envelope. For each site in
the regulatory register, the nearest building footprint is matched, the geometric volume is
computed as footprint area times eaves height, and the ratio of declared to geometric volume
is recorded. Ratios outside $[0.05, 20]$ are discarded as either an aberrant declaration or a
mismatched building. The production constant of $0.80$ is attributed in the source to a
median over 122 matched sites with an interquartile range of $[0.52, 1.22]$.

That run cannot be reproduced from the retained artifacts, and the discrepancy is recorded
here in full. The calibration script writes its detailed output to a file that is not present
in the repository, so the only auditable sample is the 68-site corpus retained for other
purposes. Recomputing the identical ratio over those 68 sites gives a median of $0.966$, an
interquartile range of $[0.80, 2.00]$, and $47\%$ of sites inside the $[0.7, 1.4]$ band,
against the $70\%$ the calibration script itself prints as its acceptance target. The
available evidence does not distinguish between a biased remnant sample and an understated
production constant. Recollecting the register cleanly, with the departmental loop and
pagination the script provides, and rederiving the factor, is the resolution.

One matching rule is worth recording because it changes the answer materially. Only the
largest building within the search radius is matched, and neighbouring buildings are not
summed: summing them moves the median ratio from $0.80$ to $0.58$, because a declared volume
belongs to one warehouse while the footprints of a whole industrial estate do not.

\subsection*{Annotation Protocol}

The dataset comprises 442 georeferenced image chips, each retained at a single location, a
single resolution and a single vintage, with a sidecar file recording its geographic
transform. Of these, 351 form the training split and 91 the validation split.

Annotators mark one class, the continuous extent of facade served by a dock apron, and the
following rules apply.

\begin{itemize}
    \item Mark an equipped facade even where no vehicle is parked. The apron and canopy
    remain visible, and a facade annotated only when a vehicle is present teaches the model
    to detect vehicles.
    \item Do not annotate a trailer holding yard that is not adjacent to a wall.
    \item Do not annotate an aeronautical tarmac or a small-business forecourt. Both are
    paved and deep but serve no docks.
    \item A facade without docks is annotated empty. That is a useful negative, not an
    omission, and roughly 40\% of the target sample is negative by design.
    \item Ignore everything above the facade line, which is roof. Rooflights and sawtooth
    roofs are perfectly periodic and are the principal source of false positives.
    \item Annotate every visible dock zone in a chip regardless of which building owns it.
    Ownership is resolved downstream by geometric clipping, so the model is never asked to
    infer it from pixels, and annotating a zone in one chip but not in the overlapping
    neighbour would give the model two contradictory labels for the same ground.
\end{itemize}

Sampling is stratified by footprint area across three bands, roughly 1\,500 to
5\,000~m\textsuperscript{2}, 5\,000 to 20\,000~m\textsuperscript{2} and above
20\,000~m\textsuperscript{2}, and across five site categories covering small enterprises,
aerospace, refrigerated, very large and retail distribution. An unstratified random draw
over-represents large warehouses, and the model inherits that bias.

\subsection*{Training Configuration}

The model is a segmentation network trained from pre-trained weights at a $640$~pixel input
size, with rotation augmentation over the full circle, horizontal and vertical flips, and
scale jitter. Rotation invariance is obtained by augmentation rather than by geometric
rectification of the facade, which is what allows the model to accept any building shape and
orientation without requiring a footprint as input. Training ran on CPU throughout, with no
GPU available, which is the practical reason the input size was not raised further and the
context for the negative resolution result of Table~\ref{tab:volume_training}.

\subsection*{Measured Facts Recorded for Reuse}

\begin{itemize}
    \item Empty bay contrast at verified locations: $0.18\sigma$, $0.22\sigma$ and
    $0.01\sigma$, against $5.3\sigma$ to $6.4\sigma$ at occupied bays on the same facades,
    over a full sweep of assumed apron depths from $0.3$~m to $20$~m.
    \item Dock pitch coherence between facades of the same building: $0.03$ to $0.05$~m,
    which provides a free quality control.
    \item Roof overhang beyond the footprint edge: $0$ to $14.4$~m. This is why the footprint
    edge is not the start of the apron, and it was the cause of an earlier defect in which
    rooflights were counted as bays.
    \item Trailer distance to building: median up to $126$~m on large platforms; at one site
    only 8 of 89 detected trailers lay within $26$~m. Trailers are not a proxy for docks.
    \item Zero-shot coverage before the learned model: 25 of 70 sites, or $36\%$.
    \item Prompt sensitivity of the promptable segmentation model used for site sampling: a
    prompt naming an industrial warehouse returns 88\,752~m\textsuperscript{2}, some 45\% of
    the view, where a prompt naming a warehouse roof returns the
    13\,478~m\textsuperscript{2} of the roof itself. The first is useful for locating a site,
    the second for measuring it, and they are not interchangeable.
\end{itemize}

\subsection*{Licensing Constraint}

The detection and segmentation library used is distributed under a strong copyleft licence.
This has no consequence for internal research use, but it becomes determinative the moment
the capability is distributed or exposed over a network. The options are a commercial licence
from the vendor or replacement of the backbone with a permissively licensed equivalent. This
is recorded as a governance decision to be taken deliberately rather than discovered at
deployment.
